\section{Introduction and Motivation}
\label{sec:intro}

\subsection{Problem}
\textit{In N\&T} (\texttt{innt.vn}) is a printing-and-packaging company in Hanoi,
Vietnam. Its catalog website presents 19 products across seven categories
(envelopes, flyers, folded leaflets, paper bags, labels/stickers, business
forms, and name cards) together with company information. Prospective B2B
customers routinely ask the same questions before requesting a quote: \emph{what
materials and sizes are available, what is the minimum order quantity, how long
does production take, and which product fits a given use case?} Answering these
manually does not scale, and a generic chatbot would either lack the
domain-specific Vietnamese printing vocabulary (e.g.\ \textit{cán màng},
\textit{bế khuôn}, \textit{giấy Couche}) or hallucinate specifications
and --- worst of all --- prices.

\subsection{Motivation for RAG}
A purely generative large language model (LLM) cannot be trusted to state
product facts: it has never seen this private catalog and will invent
plausible-but-wrong specifications. Retrieval-Augmented Generation
\cite{lewis2020rag} solves exactly this: the model answers \emph{only} from
documents retrieved from a trusted knowledge base, which (i) grounds every
answer in the real catalog, (ii) lets us update facts by editing data rather
than retraining, and (iii) makes it straightforward to enforce business rules
such as ``never quote a price --- redirect to Zalo.'' RAG is therefore the
natural fit for a small, high-precision, frequently-updated product catalog.

\subsection{Goals and contributions}
The project has two intertwined goals: ship a working Vietnamese chatbot, and
use it as an experimental platform to study \emph{how each RAG technique
contributes to quality}. Concretely, this report contributes:
\begin{enumerate}[leftmargin=1.4em,itemsep=2pt]
  \item A modular, configuration-driven RAG pipeline in which every component
        (embedding model, retrieval strategy, re-ranker, query enhancer, LLM)
        is independently swappable for controlled experiments
        (\Cref{sec:architecture}).
  \item An 80-question Vietnamese evaluation set with gold passage labels, and
        an evaluation harness combining RAGAS \cite{es2023ragas} with classical
        IR metrics (\Cref{sec:methodology}).
  \item A \emph{progressive optimization study}: starting from naive RAG and
        adding one component at a time, we quantify the marginal effect of each
        (\Cref{sec:results}).
  \item A trade-off analysis showing that the most ``advanced'' configuration is
        not the best --- query enhancement is measured and \emph{rejected} ---
        which is the core engineering lesson of the project
        (\Cref{sec:discussion}).
\end{enumerate}

\subsection{Background: three RAG architectures}
We frame the work around three progressively more capable architectures, each
building on the previous one:
\begin{description}[leftmargin=1.6em,style=nextline,itemsep=3pt]
  \item[Naive RAG (baseline).] Embed the query, perform single-stage dense
        vector search, and generate from the top-$k$ documents
        \cite{lewis2020rag}.
  \item[Advanced RAG.] Add query enhancement (HyDE \cite{gao2023hyde}), hybrid
        retrieval combining dense embeddings with sparse BM25
        \cite{robertson2009bm25}, and cross-encoder re-ranking
        \cite{nogueira2019passage}.
  \item[Agentic RAG.] Add LLM-based intent classification that routes each query
        to a specialized handler (product, comparison, business,
        recommendation, pricing, image, out-of-scope) \cite{gao2023survey,
        asai2023selfrag}.
\end{description}
Rather than only comparing these three end-to-end, we decompose them into their
constituent components and add each in isolation, which is what makes the
improvement (or regression) of every technique measurable.
